\documentclass[11pt,a4paper]{article}

\usepackage[margin=1in]{geometry}
\usepackage[T1]{fontenc}
\usepackage{lmodern}
\usepackage[english]{babel}
\usepackage{microtype}
\usepackage{booktabs}
\usepackage{tabularx}
\usepackage{array}
\usepackage{xcolor}
\usepackage{enumitem}
\usepackage{hyperref}
\usepackage{listings}
\usepackage{fancyhdr}
\usepackage{titlesec}

\hypersetup{
    colorlinks=true,
    linkcolor=blue!55!black,
    urlcolor=blue!65!black,
    pdftitle={Rover Mapping and Autonomous Exploration},
    pdfauthor={Rover Mapping Project}
}

\definecolor{codebackground}{RGB}{247,248,250}
\definecolor{codeframe}{RGB}{220,224,230}
\definecolor{codekeyword}{RGB}{35,75,145}
\definecolor{codecomment}{RGB}{80,125,90}

\lstdefinestyle{terminal}{
    backgroundcolor=\color{codebackground},
    basicstyle=\ttfamily\small,
    frame=single,
    rulecolor=\color{codeframe},
    breaklines=true,
    columns=fullflexible,
    keepspaces=true,
    showstringspaces=false
}
\lstset{style=terminal}

\pagestyle{fancy}
\fancyhf{}
\lhead{Rover Mapping Project}
\rhead{Technical Report}
\cfoot{\thepage}
\setlength{\headheight}{14pt}

\titleformat{\section}{\Large\bfseries}{\thesection}{0.7em}{}
\titleformat{\subsection}{\large\bfseries}{\thesubsection}{0.7em}{}
\setlist[itemize]{leftmargin=1.5em,itemsep=0.25em,topsep=0.4em}
\setlist[enumerate]{leftmargin=1.7em,itemsep=0.3em,topsep=0.4em}

\title{
    \vspace{-1.5em}
    \textbf{Rover Mapping and Autonomous Exploration}\\[0.4em]
    \large ROS 2 Project Technical Report
}
\author{Rover Mapping Project}
\date{\today}

\begin{document}
\maketitle
\begin{abstract}
This report describes a ROS 2 rover simulation workspace that generates a
procedural terrain, tracks rover position and battery state, publishes a local
field of view (FOV), and selects movements in exploration or destination mode.
The goal is placed inside a randomly generated target square and is disclosed
to the explorer only when the goal cell appears in an FOV message. The project
uses custom ROS 2 messages and services to connect its map, battery, and
exploration nodes.
\end{abstract}
\newpage

\section{Project Overview}

The project is a ROS 2 prototype for mapping and exploring a two-dimensional
grid world. Each map cell has four environment values:
\begin{itemize}
    \item height,
    \item roughness,
    \item muddiness, and
    \item wetness.
\end{itemize}

The simulated rover has integer grid coordinates and a floating-point battery
level. It requests movement through a ROS service. The map node checks and
charges the movement's battery cost before changing the rover position. A
local FOV message provides the explorer with nearby terrain and target-area
information.

\subsection{Workspace packages}

\begin{tabularx}{\textwidth}{>{\ttfamily}p{0.25\textwidth}X}
\toprule
\normalfont\textbf{Package} & \textbf{Responsibility}\\
\midrule
custom\_msgs & Defines the project's ROS messages and services.\\
world\_map & Implements the map, battery, and explorer nodes, launch file,
configuration, and switch-mode client.\\
\bottomrule
\end{tabularx}

\section{System Architecture}

\subsection{Nodes}

\begin{description}[style=nextline,leftmargin=1.6em,labelwidth=1.4em]
    \item[\texttt{map\_node}]
    Generates the terrain and random coordinates, owns the rover's position,
    serves movement requests, publishes the local FOV, and displays terrain
    maps.
    \item[\texttt{battery\_node}]
    Maintains battery state, publishes it periodically, and serves requests
    to deduct movement costs.
    \item[\texttt{explore\_node}]
    Consumes FOV messages, maintains an observed-cell matrix and visualization,
    and chooses moves in exploration or destination mode.
    \item[\texttt{switch\_mode\_client}]
    Makes a one-shot request to switch the explorer between its two modes
    using a boolean ROS parameter.
\end{description}

\subsection{ROS communication}

\begin{tabularx}{\textwidth}{p{0.25\textwidth}p{0.29\textwidth}X}
\toprule
\normalfont\textbf{Name} & \textbf{Type} & \textbf{Purpose}\\
\midrule
\path{/rover/fov} & Topic, \path{custom_msgs/msg/Fov} & Publishes rover pose,
battery, target-square metadata, local terrain layers, and per-cell goal
indicators.\\
\path{/rover/battery} & Topic, \path{std_msgs/msg/Float32} & Publishes current
battery state.\\
\path{/rover/move} & Service, \path{custom_msgs/srv/Move} & Requests a grid
displacement and battery deduction.\\
\path{/rover/update_battery} & Service, \path{custom_msgs/srv/Battery} &
Accepts or rejects a battery deduction.\\
\path{/rover/switch_mode} & Service, \path{custom_msgs/srv/SwitchMode} &
Changes the active explorer timer callback.\\
\bottomrule
\end{tabularx}

\section{Terrain and Random World Setup}

\subsection{Terrain generation}

The map node creates height, roughness, muddiness, and wetness grids with
dimensions \texttt{length} by \texttt{width}. Height is initialized at the map
center and expanded outward in increasing distance from that center. Each new
height value uses already-initialized neighboring values and a random noise
term. Roughness, muddiness, and wetness are independently sampled in the
interval from 0 to 1.

The map node displays four terrain plots. The height plot also shows a star
marker for the rover and updates it when an accepted move changes the rover
position.

\subsection{Coordinate generation}

The map node generates coordinates after generating the terrain. The
configurable \texttt{target\_rad} determines a square target area with side
length \(2r+1\), where \(r=\texttt{target\_rad}\). The generated values are
integers:
\begin{itemize}
    \item \textbf{Goal \((g_x,g_y)\):} sampled from map coordinates.
    \item \textbf{Target center \((t_x,t_y)\):} selected so the goal lies
    within the target square. The complete square must fit inside the map.
    \item \textbf{Spawn \((s_x,s_y)\):} selected outside the target square.
    Its octile distance to the goal must be between \(1.5r\) and \(2r\),
    inclusive.
\end{itemize}

The current generator also requires the goal and spawn cells to be dry enough
to walk on and the spawn to have at least one walkable adjacent cell. It makes
up to 10,000 attempts and raises an error if it cannot find a valid
configuration. On success, the map node logs the generated goal, target
center, radius, and spawn.

The explorer does not receive the goal coordinates directly. The FOV message
contains target-center and radius metadata, plus a boolean array indicating
whether each visible FOV cell is the goal. The explorer records the goal only
when that indicator appears in a received FOV.

\section{Messages and Services}

\subsection{Custom message: \texttt{Fov}}

The FOV message includes:
\begin{itemize}
    \item rover coordinates \texttt{x} and \texttt{y},
    \item battery level and FOV side length,
    \item target center and target radius,
    \item height, roughness, muddiness, and wetness arrays, and
    \item a boolean \texttt{goal\_cells} array.
\end{itemize}

Each terrain layer and the goal indicator are serialized as a one-dimensional,
row-major array. For FOV side length \(n\), the arrays contain \(n^2\) entries;
row \(i\), column \(j\) is at index \(i n+j\). The FOV timer publishes twice
per second. Out-of-map terrain padding uses height \(100.0\) and \(-1.0\) for
roughness, muddiness, and wetness. Out-of-map cells are not marked as goals.

\subsection{Custom services}

\begin{description}[style=nextline,leftmargin=1.7em]
    \item[\texttt{Move}]
    Request fields are \texttt{delta\_x}, \texttt{delta\_y}, and floating-point
    \texttt{deduction}. The response contains \texttt{accepted} and
    \texttt{reason}.
    \item[\texttt{Battery}]
    Accepts a floating-point \texttt{deduction} and returns \texttt{accepted}
    and \texttt{reason}. The battery node rejects a deduction that would make
    the balance negative; otherwise it subtracts the requested amount.
    \item[\texttt{SwitchMode}]
    Accepts boolean \texttt{mode} and returns \texttt{success} and
    \texttt{message}. True selects exploration mode; false selects destination
    mode.
\end{description}

\section{Movement and Battery Flow}

For each movement request, the map node first waits for the battery service.
It forwards the requested deduction to the battery node. If the battery node
accepts the deduction, the map node applies \texttt{delta\_x} and
\texttt{delta\_y} to the rover's coordinates and returns an accepted response.
If the battery service is unavailable or rejects the deduction, the rover
position is not updated.

The explorer sends movement requests asynchronously and waits for the updated
FOV before planning another move. This prevents it from planning repeatedly
against a stale rover pose.

\section{Exploration and Destination Modes}

\subsection{Exploration mode}

The explorer starts with \texttt{mode=True}. Its one-second timer calls
\texttt{moveRequest}. The planner evaluates adjacent moves using local
walkability, battery affordability, terrain values, and the number of newly
observed cells. It uses a short lookahead when scoring moves and falls back to
progress toward an unexplored frontier when the immediate information score
does not select a move. Cells covered by an incoming FOV are marked explored
in the explorer's boolean matrix.

\subsection{Destination mode}

The same one-second timer calls \texttt{destinationRequest} after the mode is
switched to false. Before the goal is observed, destination planning is toward
the target square and the planner can explore cells within that area. Once
the goal indicator is observed, the planner selects walkable, affordable
steps that reduce octile distance to the discovered goal. It returns a
zero-displacement result at the goal, and the timer callback then leaves the
rover stationary.

The planner uses local FOV data and a one-cell movement neighborhood. It is a
heuristic local planner, not a global path planner; obstacles or battery
constraints can prevent it from finding a route even when a longer route
exists.

\subsection{Mode-switch client}

The switch-mode client declares a boolean parameter named \texttt{mode}, whose
default is false. It waits up to five seconds for the switch service and for
the response, then logs the service result. The launch file currently starts
this client with \texttt{mode=False}; therefore, the launched system switches
to destination mode shortly after startup. Remove or comment out that launch
action to leave the explorer in its default exploration mode.

\section{Configuration and Running}

The default configuration is in
\texttt{src/world\_map/config/rover.params.yaml}. It currently sets a
\(100\times100\) map, FOV side length 5, target radius 5, and maximum battery
100.0. Spawn and target coordinates are generated at runtime rather than
loaded from this YAML file.

Build and launch from the workspace root:
\begin{lstlisting}
colcon build --symlink-install
source install/setup.zsh
ros2 launch world_map map.launch.py
\end{lstlisting}

To switch modes manually while the system is running, invoke the client with
the desired parameter:
\begin{lstlisting}
ros2 run world_map switch_mode_client --ros-args -p mode:=true
ros2 run world_map switch_mode_client --ros-args -p mode:=false
\end{lstlisting}

The map and explored-state visualizations use Matplotlib. A graphical
environment and an interactive Matplotlib backend are needed to see the plot
windows.

\section{Implementation Notes and Limitations}

\begin{itemize}
    \item The FOV arrays are flattened for ROS message transport; subscribers
    reshape them into \(n\times n\) matrices using the advertised FOV size.
    \item The goal is hidden as a coordinate, but the target center and radius
    are sent in every FOV message so the explorer can search the target area.
    \item The current coordinate generator checks the local terrain condition
    of the goal, spawn, and at least one spawn neighbor. It does not prove
    global connectivity between spawn, target square, and goal.
    \item The explorer's observed matrix marks cells covered by the FOV. It
    therefore represents observed cells, not only rover-occupied cells.
    \item Destination movement is greedy with respect to locally visible
    walkable cells after goal discovery. It does not perform global
    pathfinding or guarantee recovery from a locally blocked route.
    \item The map and explorer receive map dimensions through ROS parameters.
    If dimensions are changed, the values configured for both nodes must
    remain consistent.
\end{itemize}

\section{Project File Map}

\begin{tabularx}{\textwidth}{p{0.52\textwidth}X}
\toprule
\normalfont\textbf{Path} & \textbf{Role}\\
\midrule
\path{src/custom_msgs/msg/Fov.msg} & Local terrain and goal-observation message.\\
\path{src/custom_msgs/msg/Info.msg} & Rover state message.\\
\path{src/custom_msgs/srv/Move.srv} & Rover movement request and response.\\
\path{src/custom_msgs/srv/Battery.srv} & Battery deduction request and response.\\
\path{src/custom_msgs/srv/SwitchMode.srv} & Explorer mode-switch request and response.\\
\path{src/world_map/world_map/map_node.py} & Terrain generation, movement service,
FOV publisher, and map visualization.\\
\path{src/world_map/world_map/battery_node.py} & Battery publisher and deduction
service.\\
\path{src/world_map/world_map/explore_node.py} & FOV consumer, exploration planner,
destination planner, and explored-map display.\\
\path{src/world_map/world_map/switch_mode_client.py} & One-shot mode-switch client.\\
\path{src/world_map/launch/map.launch.py} & Launches the ROS nodes and configured
mode-switch client.\\
\path{src/world_map/config/rover.params.yaml} & Default map, FOV, target-radius, and
battery parameters.\\
\bottomrule
\end{tabularx}

\end{document}
